\section{Weak Interaction as Structural Fluctuation of Degenerate Closure}
\label{sec:weak}

\subsection{Reformulating the Weak Interaction}

In the Standard Model, the weak interaction is described by
an $\mathrm{SU}(2)$ gauge symmetry and couples only to left-handed fermions.
Within the conventional framework,
both the origin of $\mathrm{SU}(2)$ and the left-handed asymmetry
are postulated rather than derived.
Within the present approach, the question is reformulated:
\begin{equation}
  \boxed{
    \text{Can the weak interaction be understood as an internal
    structural process?}
  }
\end{equation}

\subsection{Degenerate Triangular Closure}

Section~\ref{sec:particles} identified degenerate triangular closure
as a reduced structural regime.
In this limit, the area approaches zero,
one degree of freedom collapses,
and two effective internal states remain:
\begin{equation}
  \boxed{|L\rangle \quad \text{and} \quad |R\rangle.}
\end{equation}
This two-state structure forms the basis for internal transitions.

\subsection{Reduction of Degrees of Freedom}

The transition from non-degenerate to degenerate closure
reduces the number of effective internal directions:
\begin{equation}
  \mathrm{SU}(3) \;\longrightarrow\; \mathrm{SU}(2).
\end{equation}
\begin{equation}
  \boxed{
    \text{Degeneracy induces a reduction of symmetry from }
    \mathrm{SU}(3) \text{ to } \mathrm{SU}(2).
  }
\end{equation}

\subsection{Internal Mixing as Structural Fluctuation}

The two-state system is not static.
It admits internal transformations that preserve reduced closure structure
while mixing the internal states:
\begin{equation}
  \begin{pmatrix}|L\rangle\\|R\rangle\end{pmatrix}
  \;\longrightarrow\;
  U\begin{pmatrix}|L\rangle\\|R\rangle\end{pmatrix},
  \qquad U \in \mathrm{SU}(2).
\end{equation}
\begin{equation}
  \boxed{
    \text{Weak interaction can be interpreted as mixing within
    degenerate closure.}
  }
\end{equation}

\subsection{Origin of Chirality}

A key feature of the weak interaction is its preference for left-handed states.
Within this framework, perfect degeneracy would imply symmetry
between $|L\rangle$ and $|R\rangle$.
However, from the foundational principle established in Vol.~1:
\begin{equation}
  \boxed{\text{Perfect symmetry is dynamically inaccessible.}}
\end{equation}
Residual asymmetry is unavoidable.
One state becomes structurally favored; the other is suppressed.
\begin{equation}
  \boxed{
    \text{Chirality can be interpreted as a consequence of residual
    asymmetry in degenerate closure.}
  }
\end{equation}

\subsection{Structural Interpretation of Decay}

Weak processes appear as particle decay.
Within the present framework,
such processes correspond to structural reorganization:
\begin{equation}
  |L\rangle \;\leftrightarrow\; |R\rangle.
\end{equation}
\begin{equation}
  \boxed{
    \text{Decay can be interpreted as reconfiguration of closure structure.}
  }
\end{equation}
This shifts the description from force-mediated interaction
to structural transformation.

\subsection{Relation to CP Asymmetry}

Since perfect symmetry is inaccessible,
residual imbalance may also manifest as asymmetry
under combined transformations.
\begin{equation}
  \boxed{
    \text{CP asymmetry can be interpreted as a structural consequence
    of incomplete degeneracy.}
  }
\end{equation}
This suggests that symmetry violations reflect underlying geometric constraints,
rather than constituting independent phenomena.

\subsection{Structural Interpretation}

\begin{equation}
  \boxed{
    \text{Weak interaction is not introduced as a force, but interpreted
    as fluctuation within a degenerate closure structure.}
  }
\end{equation}
This provides a unified structural basis for $\mathrm{SU}(2)$ symmetry,
state mixing, chirality, and decay processes.
